% =====================================================================
% SECTION TO PASTE INTO THE SHARED OVERLEAF DRAFT
% Author: Zibo.  Topic: frozen Chronos-2 layer-wise probing.
% Requires \usepackage{booktabs} in the preamble. Adjust \section level as needed.
% TODO before submit: verify author lists/venues against the arXiv PDFs you downloaded,
% and fill the manifold-paper authors (your lab's paper) in the .bib below.
% =====================================================================

\section{Intermediate-layer transferability in a native time-series foundation model}
\label{sec:chronos-probing}

\subsection{Setup}
We probe \emph{Chronos-2}, an encoder-only ($\sim$120M-parameter, 12-layer, T5-style)
time-series \emph{forecasting} foundation model, as a frozen feature extractor. For each
encoder layer $\ell \in \{0,\dots,11\}$ we read out the hidden state, mean-pool over the
content patches (each channel passed through the encoder and concatenated), and fit a single
linear probe---$\ell_2$-regularized logistic regression with features standardized on the
clean training split only, and no hyperparameter tuning. The probe's held-out accuracy
measures how \emph{linearly decodable} the class label is from layer $\ell$; a deliberately
linear probe isolates the representation rather than probe capacity, following Head2Toe
\cite{evci2022head2toe}. This mirrors the frozen-state linear-probing methodology used to
study cross-modal transfer in \cite{roger2026manifold} and intermediate-layer reuse in
\cite{roschmann2025tivit}. We evaluate on eight UEA multivariate classification tasks
spanning EEG, handwriting, gesture, chemical spectra, and astronomical light-curves; two
(Epilepsy, Cricket) saturate ($\geq 0.95$ at most layers) and are retained only as a ceiling
reference.

We test a two-part hypothesis adapted from the tunnel effect \cite{masarczyk2023tunnel}:
(i)~middle layers carry more transfer-relevant information than the final layer, and
(ii)~this middle-over-last advantage grows under distribution shift. Because Chronos-2's
final layer is specialized for the forecasting objective, it is a natural place to shed the
class-discriminative structure that intermediate layers retain.

\subsection{Result 1: intermediate layers dominate the last layer}
Across all six non-saturated tasks, probe accuracy rises through the early layers, plateaus,
and falls at the final layer; the per-layer argmax lies in the middle band (layers 3--7) in
every case. We quantify the deficit with a pre-declared statistic
$\Delta = \overline{\mathrm{acc}}(\ell\!\in\!3\text{--}8) - \mathrm{acc}(\ell\!=\!11)$
(paired bootstrap, $B{=}2000$). $\Delta$ is positive with its $95\%$ CI excluding zero in
$5/6$ tasks (Table~\ref{tab:latedrop}). The lone exception, SelfRegulationSCP2, sits at
chance at every layer ($\max_\ell|\mathrm{acc}-0.5|=0.05$) and is therefore underpowered
rather than a counterexample; LSST's effect is small but resolvable owing to its large
($2466$-sample) test set.

\begin{table}[t]
\centering
\caption{Last-layer deficit $\Delta=\overline{\mathrm{acc}}(\ell\!\in\!3\text{--}8)-\mathrm{acc}(\ell\!=\!11)$ on a frozen Chronos-2 encoder (95\% paired-bootstrap CI). Positive $\Delta$ means middle layers beat the last layer.}
\label{tab:latedrop}
\begin{tabular}{llcc}
\toprule
Dataset & Domain & $\Delta$ (95\% CI) & argmax layer \\
\midrule
UWaveGestureLibrary  & gesture        & $+0.085\ [0.051,\,0.122]$ & 4 \\
EthanolConcentration & chem.\ spectra & $+0.070\ [0.013,\,0.122]$ & 3 \\
SelfRegulationSCP1   & EEG            & $+0.063\ [0.028,\,0.098]$ & 4 \\
Handwriting          & handwriting    & $+0.050\ [0.026,\,0.075]$ & 6 \\
LSST                 & light-curves   & $+0.015\ [0.000,\,0.029]$ & 7 \\
SelfRegulationSCP2   & EEG            & $-0.018\ [-0.091,\,0.058]$ & --- \\
\bottomrule
\end{tabular}
\end{table}

This is a native-TSFM confirmation of a finding that recurs across modalities---that the most
transferable representations live in intermediate, not final, layers
\cite{evci2022head2toe, masarczyk2023tunnel, roschmann2025tivit, feofanov2026mantis}---and it
directly motivates intermediate-layer (Head2Toe-style) selection over the default last-layer
readout when reusing forecasting TSFMs for downstream discrimination.

\subsection{Result 2: OOD degradation is shift-dependent, not last-layer-concentrated}
The amplification prediction~(ii) does not hold. We apply three synthetic, label-preserving
input shifts---Gaussian noise, time-warp, and baseline drift---to the test inputs only,
score with the probe and scaler frozen on the clean training split, and measure amplification
$A = \Delta_{\mathrm{OOD}} - \Delta_{\mathrm{ID}}$. Of the $18$ non-saturated
(task $\times$ shift) cells, only $3$ are positive with CI excluding zero, against $2$ that
are significantly \emph{negative} and $13$ indistinguishable from zero---scattered across
different tasks and shifts, consistent with multiple-comparison noise ($\approx0.9$ expected
at $\alpha{=}0.05$).

The per-layer curves explain why, and reveal a more specific structure. Time-warp degrades
all layers roughly uniformly: the ID and OOD curves stay parallel (per-layer gap
std~$<0.03$), leaving no last-layer-specific penalty. Gaussian noise---the most input-like
shift---instead damages the \emph{early} layers most: on UWave and SCP1 the probe craters at
layers~0--1 and recovers by the middle (e.g.\ UWave layers~0/1 score $0.37/0.28$ against a
middle-band $0.73$), the opposite of a tunnel-style late-layer effect. This echoes
\emph{surgical fine-tuning} \cite{lee2022surgical}, where the layers that matter under shift
depend on the shift type and input-level corruptions are best addressed by adapting the
earliest layers. Our layer-wise probes give the probing-side counterpart of that observation
in a TSFM: where OOD damage localizes is shift-specific, and for input-space noise it
localizes early, not at the last layer.

\subsection{Discussion}
The robust intermediate-over-final advantage (Result~1) strengthens the case for selective,
intermediate-layer use of frozen TSFM representations. Result~2 is a deliberate negative: the
tunnel \emph{shape} is robust, but the naive corollary that synthetic input shift amplifies
it is false, and the early-layer fragility under Gaussian noise is a distinct input-coupling
phenomenon rather than a tunnel effect. Cleanly separating the late-layer deficit (tunnel
effect) from generic domain-mismatch requires a matched comparison of an \emph{in-distribution}
versus \emph{out-of-distribution} task, where the deficit should grow with distance from the
pretraining distribution if it is genuinely the tunnel effect---a natural next experiment, and
a bridge to the selective-finetuning direction.


% =====================================================================
% BIBTEX --- move these into your .bib file. VERIFY author lists/venues
% against the arXiv PDFs. Fill roger2026manifold (your lab's paper).
% =====================================================================
% @inproceedings{evci2022head2toe,
%   title={Head2Toe: Utilizing Intermediate Representations for Better Transfer Learning},
%   author={Evci, Utku and Dumoulin, Vincent and Larochelle, Hugo and Mozer, Michael C.},
%   booktitle={International Conference on Machine Learning (ICML)}, year={2022}}
%
% @inproceedings{masarczyk2023tunnel,
%   title={The Tunnel Effect: Building Data Representations in Deep Neural Networks},
%   author={Masarczyk, Wojciech and Ostaszewski, Mateusz and Imani, Ehsan and Pascanu, Razvan
%           and Mi{\l}o{\'s}, Piotr and Trzci{\'n}ski, Tomasz},
%   booktitle={Advances in Neural Information Processing Systems (NeurIPS)}, year={2023}}
%
% @article{lee2022surgical,
%   title={Surgical Fine-Tuning Improves Adaptation to Distribution Shifts},
%   author={Lee, Yoonho and Chen, Annie S. and Tajwar, Fahim and Kumar, Ananya and Yao, Huaxiu
%           and Liang, Percy and Finn, Chelsea},
%   journal={arXiv preprint arXiv:2210.11466}, year={2022}}
%
% @article{roschmann2025tivit,
%   title={Time Series Representations for Classification Lie Hidden in Pretrained Vision Transformers},
%   author={Roschmann, Simon and Bouniot, Quentin and Feofanov, Vasilii and Redko, Ievgen and Akata, Zeynep},
%   journal={arXiv preprint arXiv:2506.08641}, year={2025}}
%
% @inproceedings{feofanov2026mantis,
%   title={Mantis: Lightweight Foundation Model for Time Series Classification},
%   author={Feofanov, Vasilii and others},
%   booktitle={International Conference on Machine Learning (ICML)}, year={2026}}
%
% @article{roger2026manifold,
%   title={LLM Pretraining Shapes a Generalizable Manifold: Insights into Cross-Modal Transfer for Time Series},
%   author={FILL IN --- your lab's paper}, journal={arXiv preprint arXiv:2605.20449}, year={2026}}